%!TEX root = forallxyyc.tex
\clearpage
\thispagestyle{plain}
\onecolumn
\
\vfill

\parbox{3 in}{
Na introdução ao seu volume \emph{Symbolic Logic}, Charles Lutwidge Dodgson aconselhou: \textquotedblleft Quando chegar a qualquer passagem que não compreenda, \emph{leia-a novamente}: se \emph{ainda} não a compreender, \emph{leia-a novamente}: se não conseguir, mesmo depois de \emph{três} leituras, é muito provável que seu cérebro esteja ficando um pouco cansado. Nesse caso, guarde o livro e dedique-se a outras ocupações; no dia seguinte, quando voltar a ele descansado, é muito provável que descubra que é \emph{inteiramente} fácil.\textquotedblright{}

\medskip

O mesmo poderia ser dito deste volume, embora os leitores estejam perdoados se fizerem uma pausa para um lanche depois de \emph{duas} leituras.
}

\vfill

\parbox{3 in}{
%P.D.\ Magnus is an associate professor of philosophy in Albany, New York. His primary research is in the philosophy of science.

%\
%\\Tim Button is a University Lecturer, and Fellow of St John's College, at the University of Cambridge. His first book, \emph{The Limits of Realism}, was published by Oxford University Press in 2013.
}
\vfill
